\documentclass[10pt, letterpaper]{article}

% Packages:
\usepackage[
    ignoreheadfoot, % set margins without considering header and footer
    top=1 cm, % seperation between body and page edge from the top
    bottom=1 cm, % seperation between body and page edge from the bottom
    left=1 cm, % seperation between body and page edge from the left
    right=1 cm, % seperation between body and page edge from the right
    footskip=1.0 cm, % seperation between body and footer
    % showframe % for debugging 
]{geometry} % for adjusting page geometry
\usepackage{titlesec} % for customizing section titles
\usepackage{tabularx} % for making tables with fixed width columns
\usepackage{array} % tabularx requires this
\usepackage[dvipsnames]{xcolor} % for coloring text
\definecolor{primaryColor}{RGB}{0, 0, 0} % define primary color
\usepackage{enumitem} % for customizing lists
\usepackage{fontawesome5} % for using icons
\usepackage{amsmath} % for math
\usepackage[
    pdftitle={CV_Daniela_Jose},
    pdfauthor={Daniela José López Medina},
    pdfcreator={Overleaf},
    colorlinks=true,
    urlcolor=primaryColor
]{hyperref} % for links, metadata and bookmarks
\usepackage[pscoord]{eso-pic} % for floating text on the page
\usepackage{calc} % for calculating lengths
\usepackage{bookmark} % for bookmarks
\usepackage{lastpage} % for getting the total number of pages
\usepackage{changepage} % for one column entries (adjustwidth environment)
\usepackage{paracol} % for two and three column entries
\usepackage{ifthen} % for conditional statements
\usepackage{needspace} % for avoiding page brake right after the section title
\usepackage{iftex} % check if engine is pdflatex, xetex or luatex

% Ensure that generate pdf is machine readable/ATS parsable:
\ifPDFTeX
    \input{glyphtounicode}
    \pdfgentounicode=1
    \usepackage[T1]{fontenc}
    \usepackage[utf8]{inputenc}
    \usepackage{lmodern}
\fi
\usepackage{helvet}
\renewcommand{\familydefault}{\sfdefault}
\ttfamily

\usepackage{charter}

% Some settings:
\raggedright
\AtBeginEnvironment{adjustwidth}{\partopsep0pt} % remove space before adjustwidth environment
\pagestyle{empty} % no header or footer
\setcounter{secnumdepth}{0} % no section numbering
\setlength{\parindent}{0pt} % no indentation
\setlength{\topskip}{0pt} % no top skip
\setlength{\columnsep}{0.15cm} % set column seperation
\pagenumbering{gobble} % no page numbering

\titleformat{\section}{\needspace{4\baselineskip}\bfseries\fontsize{11pt}{11pt}\selectfont}{}{1pt}{}[\vspace{1pt}\titlerule]


\titlespacing{\section}{
    % left space:
    -1pt
}{
    % top space:
    0.3 cm
}{
    % bottom space:
    0.2 cm
} % section title spacing

\renewcommand\labelitemi{$\vcenter{\hbox{\small$\bullet$}}$} % custom bullet points
\newenvironment{highlights}{
    \begin{itemize}[
        topsep=0.10 cm,
        parsep=0.10 cm,
        partopsep=0pt,
        itemsep=0pt,
        leftmargin=0 cm + 10pt
    ]
}{
    \end{itemize}
} % new environment for highlights


\newenvironment{highlightsforbulletentries}{
    \begin{itemize}[
        topsep=0.10 cm,
        parsep=0.10 cm,
        partopsep=0pt,
        itemsep=0pt,
        leftmargin=10pt
    ]
}{
    \end{itemize}
} % new environment for highlights for bullet entries

\newenvironment{onecolentry}{
    \begin{adjustwidth}{
        0 cm + 0.00001 cm
    }{
        0 cm + 0.00001 cm
    }
}{
    \end{adjustwidth}
} % new environment for one column entries

\newenvironment{twocolentry}[2][]{
    \onecolentry
    \def\secondColumn{#2}
    \setcolumnwidth{\fill, 4.5 cm}
    \begin{paracol}{2}
}{
    \switchcolumn \raggedleft \secondColumn
    \end{paracol}
    \endonecolentry
} % new environment for two column entries

\newenvironment{threecolentry}[3][]{
    \onecolentry
    \def\thirdColumn{#3}
    \setcolumnwidth{, \fill, 4.5 cm}
    \begin{paracol}{3}
    {\raggedright #2} \switchcolumn
}{
    \switchcolumn \raggedleft \thirdColumn
    \end{paracol}
    \endonecolentry
} % new environment for three column entries

\newenvironment{header}{
    \setlength{\topsep}{0pt}\par\kern\topsep\centering\linespread{1.5}
}{
    \par\kern\topsep
} % new environment for the header

\newcommand{\placelastupdatedtext}{% \placetextbox{<horizontal pos>}{<vertical pos>}{<stuff>}
  \AddToShipoutPictureFG*{% Add <stuff> to current page foreground
    \put(
        \LenToUnit{\paperwidth-2 cm-0 cm+0.05cm},
        \LenToUnit{\paperheight-1.0 cm}
    ){\vtop{{\null}\makebox[0pt][c]{
        \small\color{gray}\textit{Last updated in November 2024}\hspace{\widthof{Last updated in November 2024}}
    }}}%
  }%
}%

% save the original href command in a new command:
\let\hrefWithoutArrow\href

\begin{document}
 \fontsize{9.5 pt}{10 pt}\selectfont 
    \newcommand{\AND}{\unskip
        \cleaders\copy\ANDbox\hskip\wd\ANDbox
        \ignorespaces
    }
    \newsavebox\ANDbox
    \sbox\ANDbox{$|$}

    \begin{header}
        \fontsize{17 pt}{17 pt}\selectfont Daniela José López Medina

        \normalsize
        \kern 5.0 pt%
        \mbox{Santiago de Surco}%
        \kern 5.0 pt%
        \AND%
        \normalsize
        \kern 5.0 pt%
        \mbox{\hrefWithoutArrow{mailto:danielalm747@gmail.com}{danielalm747@gmail.com}}%
        \kern 5.0 pt%
        \AND%
        \kern 5.0 pt%
        \mbox{\hrefWithoutArrow{tel:+51937409639}{+51 937 409 639}}%
        \kern 5.0 pt%
        \AND%
        \kern 5.0 pt%
        \mbox{\hrefWithoutArrow{https://www.linkedin.com/in/daniela-lópez/}{LinkedIn}}%

    \end{header}

    \vspace{5 pt - 0.3 cm}

    \section{Áreas de Especialización}

    \begin{onecolentry}
        \centering
        \textbf{Psicología Organizacional\kern 3.0 pt\AND\kern 3.0 pt Clima y Cultura Organizacional\kern 3.0 pt\AND\kern 3.0 pt Gestión del Potencial Humano\kern 3.0 pt\AND\kern 3.0 pt Psicoeducación}
    \end{onecolentry}

    \vspace{0.2 cm}

    \section{Educación}

        \begin{twocolentry}{
            abr 2022 - Actualidad
        }
            \textbf{Universidad de Lima}, Bachiller en Psicología\end{twocolentry}
             \vspace{0.10 cm}
        \begin{onecolentry}
            \begin{highlights}
                \item Especialización en Psicología Organizacional.
                \item Cursos destacados: Comportamiento organizacional 1, Comportamiento organizacional 2, Liderazgo, Técnicas de grupos, Recursos Humanos y Ciclo de Vida del Colaborador.
            \end{highlights}
        \end{onecolentry}

    
    \section{Proyectos destacados}
    
            \begin{twocolentry}{
            feb 2026
        }
            \textbf{Plan de Intervención en Clima y Cultura Organizacional}\end{twocolentry}

        \vspace{0.10 cm}
        \begin{onecolentry}
        Ideación de un plan de intervención para una empresa del sector logístico (bajo NDA), enfocado en la mejora continua a partir de la identificación inicial de problemáticas a través de entrevistas.
            \begin{highlights}
                \item Diseño y aplicación de un cuestionario al área de marketing para evaluar factores críticos: comunicación interna, clima laboral, cultura, diseño organizacional y potencial humano (nivel de confort y recompensas).
                \item Estructuración de un plan de intervención estimado en 8 semanas (tras 7 semanas de ideación), integrando estrategias basadas en literatura científica reciente (últimos 5 años) y bases neurológicas.
                \item Elaboración de una propuesta integral que incluyó la definición de resultados esperados, cronograma, limitaciones y presupuesto estimado (no implementado por restricciones de tiempo de la empresa).
            \end{highlights}
        \end{onecolentry}

        \vspace{0.2 cm}

        \vspace{0.2 cm}
        
        \begin{twocolentry}{
            jul 2025
        }
            \textbf{Proyecto Onboarding en Python}\end{twocolentry}

        \vspace{0.10 cm}
        \begin{onecolentry}
        El código consiste en un chatbot para automatizar atención inicial en RRHH (\hrefWithoutArrow{https://gist.github.com/Aleinad07/d74a5e9e1651bbdcbd46172c97b253a9}{Link de acceso al código}).
         \begin{highlights}
                \item Desarrollo de un código de chatbot simple que responde preguntas frecuentes sobre procesos internos de una empresa (correo, soporte, seguro médico). Incluye detección de expresiones comunes con expresiones regulares y respuestas personalizadas. 
            \end{highlights}
        \end{onecolentry}
        
        \vspace{0.2 cm}


        \begin{twocolentry}{
             2025
        }
            \textbf{Técnicas Proyectivas e Informe Psicológico}\end{twocolentry}

        \vspace{0.10 cm}
        \begin{onecolentry}
        Aplicación guiada de instrumentos proyectivos como Hombre Bajo la Lluvia, Dibujo de la Persona, Test Desiderativo y Test de Frases Incompletas de Sacks (SSCT).
            \begin{highlights}
                \item Uso de técnicas proyectivas como herramienta complementaria en procesos de análisis de perfiles y apoyo a la toma de decisiones en contextos organizacionales.
            \end{highlights}
        \end{onecolentry}

        \vspace{0.2 cm}

        \begin{twocolentry}{
            jul 2024
        }
            \textbf{Desarrollo de Material sobre Regulación Emocional y Resiliencia}\end{twocolentry}

        \vspace{0.10 cm}
        \begin{onecolentry}
        El proyecto consistió en elaborar infografías educativas para padres en comunidades vulnerables, con el objetivo de concientizar sobre la importancia de gestionar la ira en sus hijos y brindar estrategias prácticas para apoyar su desarrollo emocional.
            \begin{highlights}
                \item Diseño de infografías educativas orientadas a la regulación emocional, manejo de la ira y fortalecimiento de la resiliencia. 
                \item Identificación de señales de alerta emocionales y desarrollo de estrategias prácticas de afrontamiento.
                \item Enfoque en prevención, bienestar psicológico y promoción de un clima emocional saludable.
            \end{highlights}
        \end{onecolentry}

        \vspace{0.2 cm}

        \begin{twocolentry}{
            may 2024
        }
            \textbf{Diseño de Material Psicoeducativo para Colegio}\end{twocolentry}

        \vspace{0.10 cm}
        \begin{onecolentry}
        El proyecto consistió en la creación de un folleto psicoeducativo para el colegio Colegio San Luis Gonzaga, enfocado en fomentar la cohesión grupal y las relaciones interpersonales entre alumnos de primaria.
            \begin{highlights}
                \item Desarrollo de una charla interactiva basada en el contenido del folleto, adaptada para niños, con dinámicas que promovieron el respeto, la empatía y la colaboración.
                \item Creación de material diseñado en colaboración con la Facultad de Psicología, integrando herramientas didácticas y manualidades para reforzar los aprendizajes.
                \item Facilitación de dinámicas grupales enfocadas en empatía, respeto y trabajo en equipo.
            \end{highlights}
        \end{onecolentry}

        \vspace{0.2 cm}

        \begin{twocolentry}{
            feb 2024
        }
            \textbf{Desarrollo de Rotafolio Psicoeducativo sobre Musicoterapia}\end{twocolentry}

        \vspace{0.10 cm}
        \begin{onecolentry}
        El proyecto consistió en la creación de un rotafolio dirigido a adultos mayores, diseñado para explicarlos beneficios de la musicoterapia en aspectos físicos, cognitivos y psicosociales.
            \begin{highlights}
                \item Diseñé de contenido psicoeducativo adaptado a las necesidades de la audiencia, integrando conceptos accesibles y estrategias prácticas para fomentar el bienestar integral.
                \item Lideré sesiones educativas para promover la comprensión y aplicación de la musicoterapia como herramienta terapéutica.
            \end{highlights}
        \end{onecolentry}
        
        \vspace{0.2 cm}
        
    \section{Habilidades}

        \begin{onecolentry}
            \textbf{Tecnológicas:} Excel, Power Point, Word, Miro, Trello, Canva, Python, G Power, SQL.        
        \end{onecolentry}

        \vspace{0.2 cm}

        \begin{onecolentry}
            \textbf{Lingüísticas:} Inglés (avanzado), Español (nativo).
        \end{onecolentry}

\end{document}